\section{Conclusão}\label{sec:conclusao}

O projeto documental organiza uma CPU multiciclo para as 23 instruções da especificação 2026.2 e explicita como a extensão de pilha reutiliza o banco, a ALU e a memória. \texttt{push} conserva o operando antes de atualizar R29; \texttt{pop} lê a palavra no SP antigo antes de incrementar o ponteiro. A sequência proposta reproduz o exemplo 227 $\to$ 223 $\to$ 227 e reserva ciclos separados para a leitura síncrona e as duas escritas do \texttt{pop}.

A parte 3 disponibiliza componentes VHDL, inicialização da memória e programas Assembly que tornam mais precisa a análise das interfaces e do plano de testes. As imagens do circuito e do mapa de estados já foram incorporadas. A avaliação do funcionamento efetivo no relatório ainda depende de documentar a integração da CPU e da UC e inserir as capturas e análises das 23 instruções e das entidades construídas pela equipe. As subseções da Seção~\ref{sec:simulacoes} foram preparadas para receber esses dados sem confundir expectativa com medição. Antes da entrega, a equipe deve resolver as pendências do Apêndice~\ref{ap:pendencias} e conferir as figuras contra a implementação final.
